\subsection{紧李群与根系}

用前一节引入的术语，第 4 节与第 5 节的结果中的一个重要部分可以总结为下面的定理：

定理 2. a) (X(T), X(T/(T ∩ D(G)), R(G, T))) 是一个既约根图式；其 Weyl 群由诸 X(w)（w ∈ W）组成；群 X(C(G)) 同构于 X(T) 关于由 R(G, T) 生成的子群所得的商。

\begin{enumerate}
\def\labelenumi{\alph{enumi})}
\setcounter{enumi}{1}
\item
  \((\Gamma(\mathrm{T}),\Gamma(\mathrm{C}(\mathrm{G})_{0}),\mathrm{R}^{\vee}(\mathrm{G},\mathrm{T}))\) 是一个既约根图式；其 Weyl 群由诸 \(\Gamma(w)\)（\(w\in W\)）组成；群 \(\pi_{1}(\mathrm{G})\) 同构于 \(\Gamma(\mathrm{T})\) 关于由 \(\mathrm{R}^{\vee}(\mathrm{G},\mathrm{T})\) 生成的子群所得的商。
\item
  若把 Z-模 X(T) 与 \(\Gamma(T)\) 中的每一个都等同于另一个的对偶（第 2 节，命题 3），则前述两个根图式中的每一个都等同于另一个的逆图式。
\end{enumerate}

用 \(\mathrm{D}^{*}(\mathrm{G},\mathrm{T})\) 表示图式 \((\mathrm{X}(\mathrm{T}),\mathrm{X}(\mathrm{T}/(\mathrm{T}\cap\mathrm{D}(\mathrm{G})),\mathrm{R}(\mathrm{G},\mathrm{T})))\)，并用 \(\mathrm{D}_{*}(\mathrm{G},\mathrm{T})\) 表示图式 \((\varGamma(\mathrm{T}),\varGamma(\mathrm{C}(\mathrm{G})_{0}),\mathrm{R}^{\vee}(\mathrm{G},\mathrm{T}))\)；它们分别称为 G（关于 T）的反变图式与共变图式。

例. 1) 若 G 是秩 1 的半单群，则 \(\mathrm{D}^{*}(\mathrm{G},\mathrm{T})\) 与 \(\mathrm{D}_{*}(\mathrm{G},\mathrm{T})\) 必然同构于两个图式 \(\Delta_{2}=(\mathbf{Z},0,\{2,-2\})\)、\(\Delta_{1}=(\mathbf{Z},0,\{1,-1\})\) 之一。若 G 同构于 \(\mathbf{SU}(2,\mathbf{C})\)，则 \(\mathrm{D}_{*}(\mathrm{G},\mathrm{T})\) 同构于 \(\Delta_{1}\)（因为 G 单连通），故 \(\mathrm{D}^{*}(\mathrm{G},\mathrm{T})\) 同构于 \(\Delta_{2}\)。若 G 同构于 \(\mathbf{SO}(3,\mathbf{R})\)，则 \(\mathrm{D}^{*}(\mathrm{G},\mathrm{T})\) 同构于 \(\Delta_{1}\)（因为 \(\mathrm{C}(\mathrm{G})=\{1\})\)），故 \(\mathrm{D}_{*}(\mathrm{G},\mathrm{T})\) 同构于 \(\Delta_{2}\)。

\begin{enumerate}
\def\labelenumi{\arabic{enumi})}
\setcounter{enumi}{1}
\item
  若 G 与 \(G'\) 是两个连通紧李群，其极大环面分别为 T 与 \(T'\)，且若 \(\mathrm{D}^{*}(\mathrm{G},\mathrm{T})=(\mathrm{M},\mathrm{M}_{0},\mathrm{R})\)，\(\mathrm{D}^{*}(\mathrm{G}',\mathrm{T}')=(\mathrm{M}',\mathrm{M}_{0}',\mathrm{R}')\)，则 \(\mathrm{D}^{*}(\mathrm{G}\times\mathrm{G}',\mathrm{T}\times\mathrm{T}')\) 可等同于 \((\mathrm{M}\oplus\mathrm{M}',\mathrm{M}_{0}\oplus\mathrm{M}_{0}',\mathrm{R}\cup\mathrm{R}')\)。共变图式的情形类似。
\item
  设 N 是 T 的一个在 G 中为中心的闭子群，并设 \((M, M_{0}, R)\) 是 G 关于 T 的反变图式。则 G/N 关于 T/N 的反变图式可等同于 \((M', M_{0}', R)\)，其中 \(M'\) 是 M 中由诸 \(\lambda\)（满足 \(\lambda(N) = \{1\}\) 者）组成的子群，而 \(M_{0}' = M' \cap M_{0}\)。
\item
  类似地，设 N 是一个有限交换群，并设 \(\varphi : \pi_{1}(G) \to N\) 是一个满射同态。设 \(G'\) 是相伴于这个同态的 G 的覆盖；这是一个连通紧李群，N 是它的一个中心子群（《一般拓扑》，第十一章，编写中），且 G 可自然地等同于 \(G'/N\)。设 \(T'\) 是 \(G'\) 的作为 T 的逆像的极大环面。若 \((P, P_{0}, S)\) 是 G 关于 T 的共变图式，则 \(G'\) 关于 \(T'\) 的共变图式可等同于 \((P', P_{0}', S)\)，其中 \(P'\) 是复合同态 \(\varphi \circ f(G, T) : P \to N\) 的核（参看第 6 节，命题 11），而 \(P_{0}' = P_{0} \cap P'\)。
\end{enumerate}

注. 1) 设 \(\mathfrak{c}\) 是 \(\mathfrak{g}_{\mathbf{C}}\) 的中心；则 \(\mathfrak{c} = \mathrm{L}(\mathrm{C}(\mathrm{G}))_{(\mathbf{C})}\)。在 G 关于 T 的图式与分裂约化代数 \((\mathfrak{g}_{\mathbf{C}}, t_{\mathbf{C}})\) 的根系与逆根系之间，我们有下列关系：

\begin{enumerate}
\def\labelenumi{\alph{enumi})}
\tightlist
\item
  从 \(\mathbf{C} \otimes \Gamma(\mathrm{T})\) 到 \(\mathfrak{t}_{\mathbf{C}}\) 的典则同构诱导出从 \(\mathbf{C} \otimes \Gamma(\mathrm{C}(\mathrm{G})_0)\) 到 \(\mathfrak{c}\) 的双射，以及从 \(1 \otimes \mathrm{R}^{\vee}(\mathrm{G}, \mathrm{T})\) 到 \(2\pi i.\mathrm{R}^{\vee}(\mathfrak{g}_{\mathbf{C}}, \mathfrak{t}_{\mathbf{C}})\) 的双射。
\end{enumerate}

从 \(\mathbf{C}\otimes\varGamma(\mathrm{C}(\mathrm{G})_{0})\) 到 c 的双射，以及从 \(1\otimes\mathrm{R}^{\vee}(\mathrm{G},\mathrm{T})\) 到 \(2\pi i.\mathrm{R}^{\vee}(\mathfrak{g}_{\mathbf{C}},\mathfrak{t}_{\mathbf{C}})\) 的双射。b) 从 \(\mathbf{C}\otimes\mathrm{X}(\mathrm{T})\) 到 \(t_{C}^{*}\)（\(t_{C}\) 的对偶）的典则同构诱导出从 \(\mathbf{C}\otimes\mathrm{X}(\mathrm{T}/(\mathrm{T}\cap\mathrm{D}(\mathrm{G})))\) 到 \(t_{C}\cap\mathcal{D}(\mathfrak{g})_{C}\) 的正交补的双射，以及从 \(1\otimes\mathrm{R}(\mathrm{G},\mathrm{T})\) 到 \(\mathrm{R}(\mathfrak{g}_{\mathbf{C}},\mathfrak{t}_{\mathbf{C}})\) 的双射。

\begin{enumerate}
\def\labelenumi{\arabic{enumi})}
\setcounter{enumi}{1}
\tightlist
\item
  假设群 G 是半单的；用 R（相应地 \(R^{\vee}\)）表示根系 R(G, T)（相应地 \(R^{\vee}(G, T)\)），于是我们有包含关系
\end{enumerate}

\[
\mathrm{Q} (\mathrm{R}) \subset \mathrm{X} (\mathrm{T}) \subset \mathrm{P} (\mathrm{R}) \quad \mathrm{Q} (\mathrm{R} ^ {\vee}) \subset \Gamma (\mathrm{T}) \subset \mathrm{P} (\mathrm{R} ^ {\vee}).
\]

有限交换群 \(\mathrm{P}(\mathrm{R})/\mathrm{Q}(\mathrm{R})\) 与 \(\mathrm{P}(\mathrm{R}^{\vee})/\mathrm{Q}(\mathrm{R}^{\vee})\) 相互对偶（第六章，§1，第 9 节）；若用 \(M^{*}\) 表示有限交换群 M 的对偶群，则由前面的典则同构我们得到

\[
\begin{array}{l l} \Gamma (\mathrm{T}) / \mathrm{Q} (\mathrm{R} ^ {\vee}) \to \pi_ {1} (\mathrm{G}) & \mathrm{P} (\mathrm{R} ^ {\vee}) / \Gamma (\mathrm{T}) \to \mathrm{C} (\mathrm{G}) \\ \mathrm{P} (\mathrm{R}) / \mathrm{X} (\mathrm{T}) \to (\pi_ {1} (\mathrm{G})) ^ {\wedge} & \mathrm{X} (\mathrm{T}) / \mathrm{Q} (\mathrm{R}) \to (\mathrm{C} (\mathrm{G})) ^ {\wedge}. \end{array}
\]

特别地，\(\pi_{1}(\mathrm{G})\) 与 \(\mathrm{C}(\mathrm{G})\) 的阶的乘积等于 \(\mathrm{R}(\mathrm{G},\mathrm{T})\) 的连通指数 f（同前处）。

现在设 \(G'\) 是另一个连通紧李群，\(T'\) 是 \(G'\) 的一个极大环面。设 \(f: G \to G'\) 是一个李群同构，满足 \(f(T) = T'\)；并用 \(f_{T}\) 表示它所定义的从 T 到 \(T'\) 的同构。则 \(\mathrm{X}(f_{\mathrm{T}})\) 是从 \(\mathrm{D}^{*}(\mathrm{G}', \mathrm{T}')\) 到 \(\mathrm{D}^{*}(\mathrm{G}, \mathrm{T})\) 的同构，记作 \(\mathrm{D}^{*}(f)\)；而 \(\Gamma(f_{\mathrm{T}})\) 是从 \(\mathrm{D}_{*}(\mathrm{G}, \mathrm{T})\) 到 \(\mathrm{D}_{*}(\mathrm{G}', \mathrm{T}')\) 的同构，记作 \(\mathrm{D}_{*}(f)\)。若 \(t \in T\)，且令 \(g = f \circ \operatorname{Int} t = (\operatorname{Int} f(t)) \circ f\)，则 \(\mathrm{D}^{*}(g) = \mathrm{D}^{*}(f)\)，\(\mathrm{D}_{*}(g) = \mathrm{D}_{*}(f)\)。

命题 15. 设 \(\varphi\) 是从 \(\mathrm{D}^*(\mathrm{G}',\mathrm{T}')\) 到 \(\mathrm{D}^*(\mathrm{G},\mathrm{T})\)（相应地，从 \(\mathrm{D}_*(\mathrm{G},\mathrm{T})\) 到 \(\mathrm{D}_*(\mathrm{G}',\mathrm{T}')\)）的同构。则存在同构 \(f:\mathrm{G}\to \mathrm{G}'\)，使得 \(f(\mathrm{T}) = \mathrm{T}'\) 且 \(\varphi = \mathrm{D}^*(f)\)（相应地 \(\varphi = \mathrm{D}_*(f)\)）；若 \(f_{1}\) 与 \(f_{2}\) 是两个这样的同构，则存在元素 \(t\)（\(\mathrm{T}\) 中的），使得 \(f_{2} = f_{1}\circ \operatorname {Int}t\)。

第二个断言由第 4 节命题 9 立即得出；例如，我们对共变图式证明第一个断言。用 \(\mathfrak{g}'\)（相应地 \(\mathfrak{t}'\)）表示 \(G'\)（相应地 \(T'\)）的李代数，并用 \(\mathfrak{g}_{\mathbf{C}}'\)（相应地 \(\mathfrak{t}_{\mathbf{C}}'\)）表示其复化李代数。由第八章 §4 第 4 节定理 2 (i)，存在同构 \(\psi : \mathfrak{g}_{\mathbf{C}} \to \mathfrak{g}_{\mathbf{C}}'\)，它把 \(\mathfrak{t}_{\mathbf{C}}\) 映到 \(\mathfrak{t}_{\mathbf{C}}'\)，且在 \(\Gamma(T) \subset \mathfrak{t}_{\mathbf{C}}\) 上诱导出给定的同构 \(\varphi : \Gamma(T) \to \Gamma(T')\)。于是 \(\mathfrak{g}\) 与 \(\psi^{-1}(\mathfrak{g}')\) 是 \(\mathfrak{g}_{\mathbf{C}}\) 的两个紧形式，它们有一个公共的交 \(\mathfrak{t}\)（在 \(\mathfrak{t}_{\mathbf{C}}\) 中）；由 §3 第 2 节命题 3，存在内自同构 \(\theta\)（\(\mathfrak{g}_{\mathbf{C}}\) 的），它在 \(\mathfrak{t}_{\mathbf{C}}\) 上诱导出恒等映射，且满足 \(\theta(\mathfrak{g}) = \psi^{-1}(\mathfrak{g}')\)。把 \(\psi\) 换为 \(\psi \circ \theta\)，可以假设 \(\psi\) 把 \(\mathfrak{g}\) 映到 \(\mathfrak{g}'\)。进而，由第 2 节命题 4，存在唯一的态射 \(f_{\mathrm{T}} : T \to T'\) 使得 \(\Gamma(f_{\mathrm{T}}) = \varphi\)。于是 \(\psi\) 在 \(\mathfrak{t}\) 上的限制是 \(L(f_{\mathrm{T}})\)，且由 §2 第 6 节命题 8，存在唯一的态射 \(f : G \to G'\)，它诱导出 \(f_{\mathrm{T}}\)（在 \(T\) 上），以及 \(\psi\)（在 \(g_{C}\) 上）。把前面的论证应用于 \(\varphi^{-1}\) 与 \(\psi^{-1}\)，就得到 f 的一个逆态射，因而 f 是同构。于是 \(\mathrm{D}_{*}(f)=\Gamma(f_{\mathrm{T}})=\varphi\)，命题得证。

注意，若 T 与 \(T'\) 是 G 的两个极大环面，则图式 \(\mathrm{D}^{*}(\mathrm{G},\mathrm{T})\) 与 \(\mathrm{D}^{*}(\mathrm{G},\mathrm{T}')\) 同构（若 \(g \in G\) 满足 \(gTg^{-1} = T'\)，则 \(\operatorname{Int} g\) 是 G 到自身的同构，且把 T 映到 \(T'\)）。用 \(\mathrm{D}^{*}(\mathrm{G})\) 表示 \(\mathrm{D}^{*}(\mathrm{G},\mathrm{T})\) 的同构类（《集合论》，第 II 章，§6，第 2 节）；这是一个只依赖于 G 的根图式，称为 G 的反变图式。G 的共变图式 \(\mathrm{D}_{*}(\mathrm{G})\) 可类似地定义，于是我们得到：

推论. 两个连通紧李群 G 与 \(G'\) 同构，当且仅当图式 \(\mathrm{D}^{*}(\mathrm{G})\) 与 \(\mathrm{D}^{*}(\mathrm{G}')\)（相应地，\(\mathrm{D}_{*}(\mathrm{G})\) 与 \(\mathrm{D}_{*}(\mathrm{G}')\)）相等。

命题 16. 对每个既约根图式 D，存在一个连通紧李群 G，使得 \(\mathrm{D}^{*}(\mathrm{G})\)（相应地 \(\mathrm{D}_{*}(\mathrm{G})\)）同构于 D。

\begin{enumerate}
\def\labelenumi{\alph{enumi})}
\item
  必要时把 D 换成其逆图式，我们归结为构造 G 使 \(\mathrm{D}^{*}(\mathrm{G})\) 同构于 D。设 \(\mathrm{D} = (\mathrm{M}, \mathrm{M}_{0}, \mathrm{R})\)；则 \(Q \otimes M\) 是 \(Q \otimes M_{0}\) 与由 R 生成的向量子空间 \(\mathrm{V}(\mathrm{R})\) 的直和。此外，由于逆根在 M 上取整数值，M 在 \(\mathrm{V}(\mathrm{R})\) 上（沿 \(Q \otimes M_{0}\) 平行）的投影含于 R 的权群 \(\mathrm{P}(\mathrm{R})\)，故 M 是 \(\mathrm{M}_{0} \oplus \mathrm{P}(\mathrm{R})\) 的一个有限指数子群。用 \(D'\) 表示图式 \((\mathrm{M}_{0} \oplus \mathrm{P}(\mathrm{R}), \mathrm{M}_{0}, \mathrm{R})\)。
\item
  设 \(\mathfrak{a}\) 是一个复半单李代数，其典则根系同构于 \(R \subset \mathbf{C} \otimes V(R)\)（第八章，§4，第 3 节），并设 \(\mathfrak{g}_1\) 是 \(\mathfrak{a}\) 的一个紧实形式（\(\S 3\)，第 2 节，定理 1）。设 \(G_1\) 是一个单连通实李群，其李代数同构于 \(\mathfrak{g}_1\)；则 \(G_1\) 是紧的（\(\S 1\)，第 4 节，定理 1）。设 \(T_1\) 是 \(G_1\) 的一个极大环面。由定理 1，图式 \(D^*(G_1, T_1)\) 同构于 \((P(R), 0, R)\)。
\item
  设 \(T_{0}\) 是一个维数等于 \(M_{0}\) 的秩的环面；则 \(\mathrm{D}^{*}(\mathrm{T}_{0},\mathrm{T}_{0})\) 同构于 \((\mathrm{M}_{0},\mathrm{M}_{0},\varnothing)\)，故 \(\mathrm{D}^{*}(\mathrm{G}_{1}\times\mathrm{T}_{0},\mathrm{T}_{1}\times\mathrm{T}_{0})\) 同构于 \(D'\)（例 2）。
\item
  最后，设 N 是 \(T_{1} \times T_{0}\) 中与 M 正交的有限子群。令 \(\mathrm{G} = (\mathrm{G}_{1} \times \mathrm{T}_{0})/\mathrm{N}\)，\(\mathrm{T} = (\mathrm{T}_{1} \times \mathrm{T}_{0})/\mathrm{N}\)。则 G 是一个连通紧李群，T 是 G 的一个极大环面，且 D(G, T) 同构于 D（例 3）。
\end{enumerate}

概要. 连通紧李群按同构的分类于是归结为既约根图式的分类。连通紧半单李群对应于既约根图式 \((M, M_{0}, R)\)（满足 \(M_{0} = 0\) 者）；给出这样一个图式，等价于给出 Q 上向量空间 V 中的一个既约根系 R 以及 V 的一个满足 \(Q(R) \subset M \subset P(R)\) 的子群 M。

注 3. 设 \(\mathrm{T}'\) 是 G 的另一个极大环面，B（相应地 \(\mathrm{B}'\)）是根系 \(\mathrm{R}(\mathrm{G},\mathrm{T})\)（相应地 \(\mathrm{R}(\mathrm{G}',\mathrm{T}')\)）的一个基（第六章，§1，第 5 节，定义 2）。存在 \(g\in \mathrm{G}\) 使得 Int \(g\) 把 T 映到 \(\mathrm{T}'\) 且把 B 映到 \(\mathrm{B}'\)，且这些元素组成模 Int(T) 的唯一陪集（由于 T 与 \(\mathrm{T}'\) 共轭，可以假设 \(\mathrm{T} = \mathrm{T}'\)，于是只需应用第六章 §1 第 5 节注 4 与第 4 节命题 9）。由此可知，从 T 到 \(\mathrm{T}'\) 的同构（由 Int \(g\) 诱导）与 \(g\) 的选取无关；因而 \(\mathrm{D}_{*}(\mathrm{Int}g)\) 与 \(\mathrm{D}^{*}(\mathrm{Int}g)\) 亦然。把第八章 §5 第 3 节注 2 加以必要的改动后照搬过来，我们现在就可以定义 G 的典则极大环面、G 的典则共变与反变根图式，\ldots。

